\documentclass[11pt]{article}
\usepackage[a4paper,margin=1in]{geometry}
\usepackage{fontspec}
\usepackage[boldfont=false]{xeCJK}
\setCJKmainfont{FandolSong-Regular.otf}
\usepackage{amsmath}
\usepackage{amssymb}
\usepackage{array}
\usepackage{booktabs}
\usepackage{graphicx}
\usepackage{adjustbox}
\usepackage{enumitem}
\setlist[itemize]{topsep=2pt,partopsep=0pt,parsep=0pt,itemsep=2pt,leftmargin=1.4em}
\setlist[enumerate]{topsep=2pt,partopsep=0pt,parsep=0pt,itemsep=2pt,leftmargin=1.6em}
\usepackage{parskip}
\usepackage{fvextra}
\fvset{breaklines=true,breakanywhere=true,fontsize=\small}
\setlength{\parindent}{0pt}

\begin{document}

\begin{center}
  {\LARGE\bfseries Gas exchange in humans}\\[0.35em]
  {\large Igcse Biology}
\end{center}
\vspace{0.6em}

\section*{Gas exchange surfaces}

\textbf{Gas exchange} 气体交换 is how \textbf{oxygen} 氧气 gets into the blood and \textbf{carbon dioxide} 二氧化碳 gets out, by \textbf{diffusion} 扩散. A good gas exchange surface has four features:

\begin{itemize}
\item a large \textbf{surface area} 表面积 — so more gas can cross at once.

\item a \textbf{thin} surface — so gases have only a short distance to diffuse.

\item a good \textbf{blood supply} — to keep a steep difference in concentration.

\item good \textbf{ventilation} 通气 with air — fresh air keeps the difference steep.

\end{itemize}

\section*{The breathing system}

Air travels in and out through these parts:

\begin{center}
\adjustbox{max width=\linewidth}{%
\begin{tabular}{ll}
\toprule
\textbf{Part} & \textbf{Job} \\
\midrule
\textbf{larynx} 喉 & the voice box, at the top of the windpipe \\
\textbf{trachea} 气管 & the windpipe; carries air towards the lungs \\
\textbf{bronchi} 支气管 & two tubes, one going to each lung \\
\textbf{bronchioles} 细支气管 & smaller branching tubes inside the lungs \\
\textbf{alveoli} 肺泡 & tiny air sacs where gas exchange happens; each is wrapped in \textbf{capillaries} 毛细血管 \\
\bottomrule
\end{tabular}}
\end{center}

The alveoli make excellent gas exchange surfaces: there are millions of them (a huge surface area), each has a wall only one cell thick (a short distance), and each is surrounded by capillaries (a good blood supply).

The lungs sit in the chest, protected by the \textbf{ribs} 肋骨. Below them is a sheet of muscle, the \textbf{diaphragm} 膈肌. Between the ribs are the \textbf{intercostal muscles} 肋间肌. \textbf{(Supplement)} The trachea is held open by rings of \textbf{cartilage} 软骨, so it cannot collapse when you breathe in.

\section*{How you breathe (Supplement)}

Breathing changes the volume and \textbf{pressure} 压力 inside the chest (the \textbf{thorax} 胸腔).

\textbf{Breathing in} (\textbf{inhaling} 吸气):

\begin{itemize}
\item the external intercostal muscles contract, pulling the ribs up and out.

\item the diaphragm contracts and flattens (moves down).

\item the thorax becomes bigger, so the pressure inside drops below the outside air pressure.

\item air is pushed in.

\end{itemize}

\textbf{Breathing out} (\textbf{exhaling} 呼气) is the opposite: the muscles relax (and the internal intercostal muscles may contract), the thorax becomes smaller, the pressure rises, and air is pushed out.

\section*{Inspired and expired air}

\begin{center}
\adjustbox{max width=\linewidth}{%
\begin{tabular}{lll}
\toprule
\textbf{Gas} & \textbf{Inspired air (breathed in)} & \textbf{Expired air (breathed out)} \\
\midrule
oxygen & about 21\% & about 16\% (less) \\
carbon dioxide & about 0.04\% & about 4\% (more) \\
\textbf{water vapour} 水蒸气 & a little & a lot (more) \\
\bottomrule
\end{tabular}}
\end{center}

You can test for carbon dioxide with \textbf{limewater} 石灰水, which turns cloudy. Expired air turns limewater cloudy much faster than inspired air, showing it contains more carbon dioxide.

\section*{Breathing and exercise}

During exercise your muscles respire faster and make more carbon dioxide. This raises the \textbf{concentration} 浓度 of carbon dioxide in the blood. \textbf{(Supplement)} Your \textbf{brain} 大脑 detects the rise and makes you breathe at a faster \textbf{rate} 速率 and a greater depth. This brings in more oxygen and removes the extra carbon dioxide quickly.

\section*{Keeping the airways clean (Supplement)}

The airways are lined with two kinds of cell that trap and remove dirt:

\begin{itemize}
\item \textbf{goblet cells} 杯状细胞 make \textbf{mucus} 黏液, which traps \textbf{pathogens} 病原体 and dust \textbf{particles} 粒子.

\item \textbf{ciliated cells} 纤毛细胞 have tiny hairs that sweep the mucus, with the trapped dirt, up to the throat, where it is swallowed.

\end{itemize}

\section*{Exam tips}

\begin{itemize}
\item A gas exchange surface is large, thin, with a good blood supply and good ventilation.

\item Air path: larynx → trachea → bronchi → bronchioles → alveoli.

\item Breathing in: external intercostal muscles and diaphragm contract → thorax bigger → pressure lower → air in. Breathing out is the opposite.

\item Expired air has \textbf{less} oxygen, \textbf{more} carbon dioxide and \textbf{more} water vapour. Limewater tests for carbon dioxide.

\item Exercise → more carbon dioxide in the blood → brain → faster, deeper breathing.

\end{itemize}


\end{document}
